% Template do Relatório Parcial e Final para o Programa Iniciação Científica da UFPR, baseado em  https://www.overleaf.com/latex/templates/modelo-pibic-ufc-itapaje/jmnnwfkhgtbw.
% Versão 18/09/2023.
\documentclass[a4paper, 12pt]{article}
\usepackage[utf8]{inputenc}
\usepackage[portuguese]{babel}
\usepackage[breaklinks=true]{hyperref}
\usepackage{breakcites}
\usepackage{ragged2e}
\usepackage{mathptmx}
\usepackage{enumitem}
\usepackage{xcolor}

% Margins
\usepackage{geometry}
\geometry{left=30mm,right=20mm,bindingoffset=0mm, top=30mm,bottom=20mm,includefoot}

% Images
\usepackage{graphicx}
\graphicspath{ {./images/} }

% Title font size
\usepackage{titlesec}
\titleformat*{\section}{\bfseries}

% Indent first sentence of a paragraph
\usepackage{indentfirst}

% Page numbers at bottom right
\usepackage{fancyhdr}
\pagestyle{fancy}
\fancyhf{}
\renewcommand{\headrulewidth}{0pt}
\rfoot{\thepage}

% 1.5 line spacing (yes, 1.25 is actually 1.5)
\fontsize{10}{12}
\linespread{1.50}
\fontfamily{ptm}
\selectfont

% Fix vspacing for section titles
\let\oldincludegraphics\includegraphics
%\renewcommand{\includegraphics}[1]
%{{\vspace{20pt}\includegraphics[#1]\vspace{-10pt}}}

% Fix vspacing for section titles
% \let\oldsection\section
% \renewcommand{\section}[1]{{\vspace{20pt}\oldsection{#1}\vspace{-10pt}}}

% \let\oldsubsection\subsection
% \renewcommand{\subsection}[1]{{\vspace{20pt}\oldsubsection{#1}\vspace{-10pt}}}

\titleformat{\section}{\bfseries\uppercase}{\thesection}{1em}{}
\titleformat{\subsection}{\bfseries\uppercase}{\thesubsection}{1em}{}

% Indent size
\setlength\parindent{10mm}

% Jump 20pt
\newcommand{\jump}[1]{{\vspace{20pt}}}

\newenvironment{leftquote}[1][]
{
    \hspace{30pt} #1
    \vspace{20pt}
    \fontsize{10}{10}\selectfont
    \begin{flushright}
    \begin{minipage}{120mm}
}
{
    \end{minipage}
    \end{flushright}
    \par
    \vspace{20pt}
}

\newcommand{\tmp}{}

\newenvironment{container}[3]
{
    \renewcommand{\tmp}{#3} %gambiarra
    
    \begin{center}
    \begin{minipage}{#1}
    \centering
    #2
}
{
    \raggedright
    {\footnotesize \tmp}
    \end{minipage}
    \end{center}
}

\begin{document}

% Remove page number
\clearpage
\thispagestyle{empty}

\begin{bfseries}
\begin{center}

\vspace{0.5cm} 
UNIVERSIDADE FEDERAL DO PARANÁ \\

\vspace{0.6cm}
NOME DO ALUNO\\
\vspace{2.5cm}
RELATÓRIO FINAL \\

\vspace{1cm}

\begin{minipage}[t]{0.45\textwidth}\hspace{3cm}
\raggedright
\textbf{PROGRAMA DE IC:}\\ \hspace{3cm}
\textnormal{(X) PIBIC}\\ \hspace{3cm}
\textnormal{( ) PIBIC Af}\\ \hspace{3cm}
\textnormal{( ) PIBIC EM}\\ \hspace{3cm}
\textnormal{( ) PIBITI}\\ \hspace{3cm}

\end{minipage}
\begin{minipage}[t]{0.45\textwidth}\hspace{1cm}
\raggedright
\textbf{MODALIDADE:}\\ \hspace{1cm}
\textnormal{( ) CNPq}\\ \hspace{1cm}
\textnormal{( ) UFPR TN}\\ \hspace{1cm}
\textnormal{(X) Fundação Araucária}\\ \hspace{1cm}
\textnormal{( ) Voluntária}\\ \hspace{1cm}

\end{minipage}

\vspace*{1cm}
\textbf{\MakeUppercase{Título do Plano de Trabalho}} \\

\begin{minipage}[t]{0.45\textwidth} \hspace{2cm}
\raggedright

\end{minipage}
\begin{minipage}[t]{0.5\textwidth}\hspace{2cm}
\raggedright

\begin{flushleft}
\footnotesize  
\justify
\hspace{-0,35cm}
\textnormal{
Relatório apresentado à Coordenação de Iniciação Científica e Tecnológica da Universidade Federal do Paraná como requisito parcial da conclusão das atividades de Iniciação Científica ou Iniciação em desenvolvimento tecnológico e Inovação - Edital 20??.\\
Orientador(a): Prof.(a).  Nome Professor(a)\\
Título do Projeto: Tíulo do Projeto.
}
\end{flushleft}
    
\end{minipage}

\vspace*{3cm}
\textbf{\MakeUppercase{CURITIBA}} \\
\textbf{\MakeUppercase{\textnormal{2023}}}
\end{center}

\end{bfseries}
\newpage

\hspace{-1,1cm}
\textbf{\MakeUppercase{RESUMO}} \\ 
Deve ter no máximo 25 linhas, conter conceito do projeto (opcional), objetivo(s), justificativa ou evidências de interesse, materiais e métodos/metodologia, principais resultados e conclusão/considerações finais.

\vspace{1cm}
\hspace{-1cm}
\textbf{Palavras Chave:} Modelo Latex; Relatório Final; IC-UFPR.


\section{INTRODUÇÃO} 

Introdução do tema; justificativa; questão norteadora ou pergunta de pesquisa (se houver); hipóteses ou pressupostos (se houver); objetivo(s).


\section{REVISÃO DA LITERATURA}


\section{MATERIAIS E MÉTODOS}

Deverá conter o método de pesquisa/tipo de estudo; local de desenvolvimento ou da coleta de dados/informações; técnicas-procedimentos de coleta/obtenção dos dados; amostra/participantes; técnica-procedimentos de análise/tratamento dos dado.


\section{RESULTADOS e DISCUSSÃO}

Poderão ser redigidos juntos ou em tópicos separados.


\section{CONCLUSÕES/CONSIDERAÇÕES}


\section{REFERÊNCIAS}

Listar as referências utilizadas no texto, segundo normas da ABNT ou a utilizada na área da CAPES.

\bibliographystyle{abntex2-alf}
\bibliography{bibfile}

\end{document}
